% Чтобы глава начиналась с новой страницы, используем \include
% Чтобы текст продолжался без разрывов, используем \input и команду \let\clearpage\relax
%\begingroup
%  \let\clearpage\relax                        % Чтобы текст продолжался без разрывов
%  \par\vspace{3ex}                            % Отступ
%  \chapter{Оформление различных элементов}\label{ch:ch1}
%\endgroup
\noskipchapter{Оформление различных элементов}\label{ch:ch1}


\textit{Основное содержание посвящается раскрытию предмета исследования.
Количество глав и их характер зависят от особенностей темы, цели и задач, выбранных методов и т.\,п.
Содержание глав ВКР должно продемонстрировать умение автора логично, последовательно, сжато и аргументированно излагать материал, выявлять новое и оригинальное в ходе разработки рассматриваемой проблемы.}

\vspace{2ex}

\section{Форматирование текста}\label{sec:ch1/sec1}

Мы можем сделать \textbf{жирный текст} и \textit{курсив}.

\section{Ссылки}\label{sec:ch1/sec2}

Сошлёмся на библиографию.
Одна ссылка: \cite[с.~54]{Sokolov}.
Две ссылки: \cite{Sokolov, Gaidaenko}.
Много ссылок:
\cite{Gost.7.0.53, Razumovski, Lagkueva, Pokrovski, Methodology, Berestova, Kriger}, \cite{Sirotko, Lukina, Encyclopedia, Nasirova}.
И~ещё немного ссылок:
\cite{Article,Book,Booklet,Conference,Inbook,Incollection,Manual,Mastersthesis,Misc,Phdthesis,Proceedings,Techreport,Unpublished}.


Ссылки на собственные работы:~\cite{bib2, bib3}.

Ссылки на патент и свидетельство о регистрации программы:~\cite{patbib1, progbib1}.

Ссылка на изображение: рисунок~\cref{fig:knuth}.

Ссылка на формулу: формула~\cref{eq:equation1}.

Ссылки на приложения: приложение~\cref{app:A}.

Стандартной практикой является добавление к ссылкам префикса, характеризующего тип элемента.
Это не является строгим требованием, но~позволяет лучше ориентироваться в документах большого размера.
Например, для ссылок на~рисунки используется префикс \textit{fig}, для ссылки на~таблицу "--- \textit{tab}.

В таблице \cref{tab:tab_pref} приложения~\cref{app:A} приведён список рекомендуемых к использованию стандартных префиксов.

\section{Формулы}\label{sec:ch1/sec3}

Для добавления формул можно использовать пары
\verb+$+\dots\verb+$+ и \verb+$$+\dots\verb+$$+,
но~они считаются устаревшими.
Лучше использовать их функциональные аналоги
\verb+\(+\dots\verb+\)+ и \verb+\[+\dots\verb+\]+.

Благодаря пакету \verb|icomma|, \LaTeX~одинаково хорошо воспринимает
в~качестве десятичного разделителя и запятую (\(3,1415\)), и точку (\(3.1415\)).

\subsection{Ненумерованные одиночные формулы}\label{subsec:ch1/sec3/sub1}

Вот так может выглядеть формула, которую необходимо вставить в~строку
по~тексту: \(x \approx \sin x\) при \(x \to 0\).

А вот так выглядит ненумерованная отдельностоящая формула c подстрочными
и надстрочными индексами:
\[
    (x_1+x_2)^2 = x_1^2 + 2 x_1 x_2 + x_2^2
\]

Формула с неопределенным интегралом:
\[
    \int f(\alpha+x)=\sum\beta
\]

При использовании дробей формулы могут получаться очень высокие:
\[
    \frac{1}{\sqrt{2}+
        \displaystyle\frac{1}{\sqrt{2}+
            \displaystyle\frac{1}{\sqrt{2}+\cdots}}}
\]

\subsection{Ненумерованные многострочные формулы}\label{subsec:ch1/sec3/sub2}

Вот так можно написать две формулы, не нумеруя их, чтобы знаки <<равно>> были
строго друг под другом:
\begin{align}
    f_W & =  \min \left( 1, \max \left( 0, \frac{W_{soil} / W_{max}}{W_{crit}} \right)  \right), \nonumber \\
    f_T & =  \min \left( 1, \max \left( 0, \frac{T_s / T_{melt}}{T_{crit}} \right)  \right), \nonumber
\end{align}

Выровнять систему ещё и по переменной \( x \) можно, используя окружение
\verb|alignedat| из пакета \verb|amsmath|. Вот так:
\[
|x| = \left\{
\begin{alignedat}{2}
    &&x, \quad &\text{eсли } x\geqslant 0 \\
    &-&x, \quad & \text{eсли } x<0
\end{alignedat}
\right.
\]
Здесь первый амперсанд (в исходном \LaTeX\ описании формулы) означает
выравнивание по~левому краю, второй "--- по~\( x \), а~третий "--- по~слову
<<если>>. Команда \verb|\quad| делает большой горизонтальный пробел.

Ещё вариант:
\[
    |x|=
    \begin{cases}
        \phantom{-}x, \text{если } x \geqslant 0 \\
        -x, \text{если } x<0
    \end{cases}
\]

Кроме того, для  нумерованных формул \verb|alignedat| делает вертикальное
выравнивание номера формулы по центру формулы. Например, выравнивание
компонент вектора:
\begin{equation}
    \label{eq:2p3}
    \begin{alignedat}{2}
        {\mathbf{N}}_{o1n}^{(j)} = \,{\sin} \phi\,n\!\left(n+1\right)
        {\sin}\theta\,
        \pi_n\!\left({\cos} \theta\right)
        \frac{
        z_n^{(j)}\!\left( \rho \right)
        }{\rho}\,
        &{\boldsymbol{\hat{\mathrm e}}}_{r}\,+   \\
        +\,
        {\sin} \phi\,
        \tau_n\!\left({\cos} \theta\right)
        \frac{
        \left[\rho z_n^{(j)}\!\left( \rho \right)\right]^{\prime}
        }{\rho}\,
        &{\boldsymbol{\hat{\mathrm e}}}_{\theta}\,+   \\
        +\,
        {\cos} \phi\,
        \pi_n\!\left({\cos} \theta\right)
        \frac{
        \left[\rho z_n^{(j)}\!\left( \rho \right)\right]^{\prime}
        }{\rho}\,
        &{\boldsymbol{\hat{\mathrm e}}}_{\phi}\:.
    \end{alignedat}
\end{equation}

Ещё об отступах.
Иногда для лучшей <<читаемости>> формул полезно немного исправить стандартные интервалы \LaTeX\ с учётом логической структуры самой формулы.
Например в формуле~\cref{eq:2p3} добавлен небольшой отступ \verb+\,+ между основными сомножителями, ниже результат применения всех вариантов отступа:
\begin{align*}
    \backslash!             & \quad f(x) = x^2\! +3x\! +2         \\
    \mbox{по-умолчанию}     & \quad f(x) = x^2+3x+2               \\
    \backslash,             & \quad f(x) = x^2\, +3x\, +2         \\
    \backslash{:}           & \quad f(x) = x^2\: +3x\: +2         \\
    \backslash;             & \quad f(x) = x^2\; +3x\; +2         \\
    \backslash \mbox{space} & \quad f(x) = x^2\ +3x\ +2           \\
    \backslash \mbox{quad}  & \quad f(x) = x^2\quad +3x\quad +2   \\
    \backslash \mbox{qquad} & \quad f(x) = x^2\qquad +3x\qquad +2
\end{align*}

Посмотрим на систему уравнений на примере аттрактора Лоренца:

\[
\left\{
\begin{array}{rl}
    \dot x = & \sigma (y-x)  \\
    \dot y = & x (r - z) - y \\
    \dot z = & xy - bz
\end{array}
\right.
\]

А для вёрстки матриц удобно использовать многоточия:
\[
    \left(
        \begin{array}{ccc}
            a_{11} & \ldots & a_{1n} \\
            \vdots & \ddots & \vdots \\
            a_{n1} & \ldots & a_{nn} \\
        \end{array}
    \right)
\]

\subsection{Нумерованные формулы}\label{subsec:ch1/sec3/sub3}

А вот так пишется нумерованная формула:
\begin{equation}
    \label{eq:equation1}
    e = \lim_{n \to \infty} \left( 1+\frac{1}{n} \right) ^n
\end{equation}

Нумерованных формул может быть несколько:
\begin{equation}
    \label{eq:equation2}
    \lim_{n \to \infty} \sum_{k=1}^n \frac{1}{k^2} = \frac{\pi^2}{6}
\end{equation}

Впоследствии на формулы~\cref{eq:equation1, eq:equation2} можно ссылаться.

Сделать так, чтобы номер формулы стоял напротив средней строки, можно,
используя окружение \verb|multlined| (пакет \verb|mathtools|) вместо
\verb|multline| внутри окружения \verb|equation|. Вот так:
\begin{equation} % \tag{S} % tag - вписывает свой текст
    \label{eq:equation3}
    \begin{multlined}
        1+ 2+3+4+5+6+7+\dots + \\
        + 50+51+52+53+54+55+56+57 + \dots + \\
        + 96+97+98+99+100=5050
    \end{multlined}
\end{equation}

Уравнения~\cref{eq:subeq_1,eq:subeq_2} демонстрируют возможности
окружения \verb|\subequations|.
\begin{subequations}
    \label{eq:subeq_1}
    \begin{gather}
        y = x^2 + 1 \label{eq:subeq_1-1} \\
        y = 2 x^2 - x + 1 \label{eq:subeq_1-2}
    \end{gather}
\end{subequations}

Ссылки на отдельные уравнения~\cref{eq:subeq_1-1,eq:subeq_1-2,eq:subeq_2-1}.
\begin{subequations}
    \label{eq:subeq_2}
    \begin{align}
        y & = x^3 + x^2 + x + 1 \label{eq:subeq_2-1} \\
        y & = x^2
    \end{align}
\end{subequations}

\subsection{Форматирование букв}

В формулах можно использовать разные математические алфавиты:
\begin{align}
    \mathcal{ABCDEFGHIJKLMNOPQRSTUVWXYZ} \nonumber  \\
    \mathfrak{ABCDEFGHIJKLMNOPQRSTUVWXYZ} \nonumber \\
    \mathbb{ABCDEFGHIJKLMNOPQRSTUVWXYZ} \nonumber
\end{align}

В формулах можно использовать греческие буквы:
\[
    \alpha\beta\gamma\delta\epsilon\zeta\eta\theta%
    \iota\kappa\lambda\mu\nu\xi\pi\rho%
    \sigma\tau\upsilon\phi\chi\psi\omega\Gamma\Delta%
    \Theta\Lambda\Xi\Pi\Sigma\Upsilon\Phi\Psi\Omega.
\]

Для жирного начертания символов в формулах используется команда \verb|\symbf|, для курсива --- \verb|\symit|:
\[%https://texfaq.org/FAQ-boldgreek
\symbf{
    \alpha\beta\gamma\delta\epsilon\zeta\eta\theta%
    \iota\kappa\lambda\mu\nu\xi\pi\rho%
    \sigma\tau\upsilon\phi\chi\psi\omega\Gamma\Delta%
    \Theta\Lambda\Xi\Pi\Sigma\Upsilon\Phi\Psi\Omega
},
\]
\[
\symit{
    \alpha\beta\gamma\delta\epsilon\zeta\eta\theta%
    \iota\kappa\lambda\mu\nu\xi\pi\rho%
    \sigma\tau\upsilon\phi\chi\psi\omega\Gamma\Delta%
    \Theta\Lambda\Xi\Pi\Sigma\Upsilon\Phi\Psi\Omega
}.
\]

За прямое и наклонное начертание в пакете \verb|unicode-math| отвечает опция \verb|math-style|.
Значение \verb|math-style=french| даёт следующий результат для латинских и греческих букв:
\begin{enumerate}
    \item Латинские буквы
    \begin{itemize}
        \item курсив для строчных букв: \(xyz\), прямое начертание для заглавных: \(XYZ\);
        \item жирный шрифт даёт прямое начертание как строчных, так и заглавных букв: \(\symbf{xyz}\), \(\symbf{XYZ}\);
    \end{itemize}
    \item Греческие буквы
    \begin{itemize}
        \item прямое начертание обычным шрифтом: \(\phi\psi\omega\), \(\Phi\Psi\Omega\);
        \item прямое начертание жирным шрифтом: \(\symbf{\phi\psi\omega}\), \(\symbf{\Phi\Psi\Omega}\).
    \end{itemize}
\end{enumerate}

Для непосредственного управления шрифтом в формулах можно использовать команды \verb|\sym*|:
\begin{itemize}
    \item \verb|\symup| --- прямое начертание: \(\symup{xyz}\), \(\symup{\phi\psi\omega}\);
    \item \verb|\symbfup| --- жирное прямое начертание: \(\symbfup{xyz}\), \(\symbfup{\phi\psi\omega}\);
    \item \verb|\symit| --- курсив: \(\symit{xyz}\), \(\symit{\phi\psi\omega}\);
    \item \verb|\symbfit| --- жирный курсив: \(\symbfit{xyz}\), \(\symbfit{\phi\psi\omega}\).
\end{itemize}

\subsection{Форматирование чисел}\label{sec:units}

Числа можно форматировать при помощи команды \verb|\num|:
\num{5.3};
\num{2.3e8};
\num{12345,67890};
\num{2,6 d4};
\num{1+-2i};
\num{.3e45};
\num[exponent-base=2]{5 e64};
\num{1.654 x 2.34 x 3.430}
\num{12 x 3/4}.
Для написания последовательности чисел можно использовать команды \verb|\numlist| и \verb|\numrange|:
\numlist{10;30;50;70}; \numrange{10}{30}.
Значения углов можно форматировать при помощи команды \verb|\ang|:
\ang{2.67};
\ang{30,3};
\ang{-1;;};
\ang{;-2;};
\ang{;;-3};
\ang{300;10;1}.

%Обратите внимание, что ГОСТ запрещает использование знака <<->> для обозначения отрицательных чисел
%за исключением формул, таблиц и~рисунков.
%Вместо него следует использовать слово <<минус>>.

С дополнительными опциями форматирования можно ознакомиться в~описании пакета \verb|siunitx|.


\section{Рецензирование текста}\label{sec:markup}

В шаблоне заданы команды рецензирования.
Они видны при компиляции шаблона в режиме черновика или при установке
соответствующей настройки (\verb+showmarkup+) в~файле \verb+common/setup.tex+.

Команда \verb+\todo+ отмечает текст красным цветом.
\todo{Например, так.}

Команда \verb+\note+ позволяет выбрать цвет текста.
\note{Чёрный, } \note[red]{красный, } \note[green]{зелёный, } \note[blue]{синий.}

Окружение \verb+commentbox+ также позволяет выбрать цвет.

\begin{commentbox}[red]
    Красный текст.

    Несколько параграфов красного текста.
\end{commentbox}

\begin{commentbox}[blue]
    Синяя формула.

    \begin{equation}
        \alpha + \beta = \gamma
    \end{equation}
\end{commentbox}

Окружение \verb+commentbox+ позволяет закомментировать участок кода в~режиме чистовика.
Чтобы убрать кусок кода для всех режимов, можно использовать окружение \verb+comment+.

\begin{comment}
    Этот текст всегда скрыт.
\end{comment}

\FloatBarrier
